%%%%%%%%%%%%%%%%%%%%%%%%%%%%%%%%%%%%%%%%
%
% $ Autor: Gruppe 07 $
% $ Projekt: Magic Faucet Fountain $
% $ Pfad: LaTeXTemplate/Nano33BLESense/Contents/General/ArduinoIDESetupNano33.tex $
% $ Kurzbeschreibung: Dieses Dokument beschreibt die Einrichtung des Arduino Nano 33 BLE Sense und zeigt am Ende ein Beispiel-Sketch zum testen des Programms. $
%
% !TeX spellcheck = de_DE/GB
% !TeX program = pdflatex
% !BIB program = biber/bibtex
% !TeX encoding = utf8
%
%%%%%%%%%%%%%%%%%%%%%%%%%%%%%%%%%%%%%%%%



\chapter{Einrichtung des Arduino Nano 33 BLE Sense}\label{ArduinoNanoBLE}

\section{Einführung}

In diesem Kapitel wird der Prozess der Verbindung des Arduino Nano 33 BLE Sense mit einem Computer oder Laptop sowie die Konfiguration der erforderlichen Grundeinstellungen beschrieben. Die folgenden Schritte werden detailliert erläutert:

\begin{itemize}
	\item Installation der notwendigen Software-Tools
	\item Korrekte Einrichtung des Boards in der Arduino IDE
	\item Verifizierung der Funktionsfähigkeit durch ein Beispielsketch
\end{itemize}

\section{Konfiguration des Arduino Nano 33 BLE Sense}

Um den \textbf{Arduino Nano 33 BLE Sense} in einer Offline-Umgebung zu programmieren, gehen Sie wie folgt vor:

\begin{enumerate}
	\subsection*{Installation des Treibers}
	\item \textbf{Installation der Arduino IDE:} Installieren Sie zunächst die neueste Version der Arduino IDE auf Ihrem Computer.
	
	\item \textbf{Installation der Pakete:} Nach Abschluss der IDE-Installation öffnen Sie die \menu{Arduino IDE} und navigieren zum Menü \menu{Werkzeuge} in der oberen linken Ecke.
	
	\item Wählen Sie im Menü \menu{Werkzeuge} den \menu{Board-Verwalter} aus.
	
	\item Im Fenster \menu{Board-Verwalter} verwenden Sie die Suchleiste, um den \textbf{Arduino Nano 33 BLE Sense} zu finden, indem Sie den Namen eingeben (siehe Abbildung \ref{fig:arduinoMbed}).
	
	\item Wählen Sie aus den Suchergebnissen \menu{Arduino Mbed OS Nano Boards} aus und klicken Sie auf \menu{Installieren}, um das benötigte Paket zu installieren.
	
\end{enumerate}

Das Mbed OS Nano Board-Paket beinhaltet die Unterstützung für mehrere Boards der Arduino Nano Familie, einschließlich des Arduino Nano 33 BLE Sense. Nach Abschluss der Installation verbinden Sie den Arduino Nano 33 BLE Sense über ein USB-A zu USB-Micro Kabel mit Ihrem Computer, um mit der Programmierung zu beginnen.

\begin{figure}[H]
	\centering
	\caption{Mbed OS Nano-Paket installation}
	\label{fig:arduinoMbed}
	\includegraphics[width=0.8\textwidth]{Arduino/ArduinoIDE/arduinoMbed.png}
\end{figure}


\subsection{Verbindung und Konfiguration des Arduino-Boards mit dem Computer}

Um entweder ein integriertes Beispiel oder ein eigenes Programm auf dem Arduino Nano 33 BLE Sense auszuführen, gehen Sie wie folgt vor:

\begin{enumerate}
	\item Verbinden Sie das Arduino-Board über ein USB-A-zu-USB-Micro-Kabel mit dem Computer.
	
	\item Öffnen Sie die Arduino IDE auf dem Computer. Es erscheint eine leere Arbeitsumgebung mit den Standardfunktionen \PYTHON{void setup()} und \PYTHON{void loop()}.
	
	\item Navigieren Sie zum Menü \menu{Werkzeuge}, wählen Sie \menu{Board} und selektieren Sie das angeschlossene Board \textbf{Arduino Nano 33 BLE Sense} (siehe Abbildung~\ref{fig:arduinoConnect}).
	
	\item Wählen Sie anschließend den korrekten Port unter \menu{Werkzeuge > Port} aus.
	
\end{enumerate}

Diese Einrichtung ermöglicht der Arduino IDE, das Arduino Nano 33 BLE Sense zu erkennen und mit ihm zu kommunizieren.

\begin{figure}[H]
	\centering
	\caption{Verbinden des Mbed OS Nano-Pakets mit dem Board}
	\label{fig:arduinoConnect}
	\includegraphics[width=0.8\textwidth]{Arduino/ArduinoIDE/arduinoConnect.png}
\end{figure}


\subsection{Auswahl des richtigen Ports}

Nach der Auswahl des Arduino Nano 33 BLE Sense Boards muss der korrekte Kommunikationsport eingestellt werden. Gehen Sie hierfür wie folgt vor:
\textbf{Bootloader-Modus aktivieren:}
\begin{enumerate}
	\item Drücken Sie den weißen Reset-Knopf auf dem Board (siehe Abbildung~\ref{fig:ArduinoNano33BLESenseResetButton})
	\item Der Knopf befindet sich in der Nähe des USB-Anschlusses
	\item Die Onboard-LED blinkt während des Bootloader-Modus
\end{enumerate}
\textbf{Port auswählen:}
\begin{enumerate}
	\item Navigieren Sie in der Arduino IDE zu \menu{Werkzeuge > Port}
	\item Wählen Sie den Port mit der Bezeichnung \menu{Arduino Nano 33 BLE Sense} aus
	\item Unter Windows erscheint dieser typischerweise als COM-Port (z.B. COM3)
\end{enumerate}

\begin{center}
	\begin{tikzpicture}
		\node (G) at (0,0) {\includegraphics[width=8cm]{Arduino/ArdiunoIDE2/resetbutton.png}};
		
		\draw[green,domain=0:360,line width=1.4pt] plot ({-1.5+0.4*cos(\x)},{-1.6
			+0.4*sin(\x)});
	\end{tikzpicture}
		\captionof{figure}{Arduino Nano 33 BLE Sense Reset Button}\label{fig:ArduinoNano33BLESenseResetButton}		
\end{center}
		

\begin{center}
	\begin{tikzpicture}
		\node (G) at (0,0) {\includegraphics[width=8cm]{Arduino/ArdiunoIDE2/resetbutton.png}};		
		\draw[green,domain=0:360,line width=1.4pt] plot ({-2.8+0.4*cos(\x)},{-0.8
			+0.4*sin(\x)});			
		\draw[orange,domain=0:360,line width=1.4pt] plot ({-2.8+0.4*cos(\x)},{-2.6
	+0.4*sin(\x)});			
	\end{tikzpicture}
		\captionof{figure}{green and orange Builtin-LED}\label{fig:ArduinoNano33BLESenseResetButton}	
\end{center}	


Nach erfolgreicher Durchführung der vorherigen Schritte muss der korrekte Port für den Upload des Programms ausgewählt werden:

\begin{enumerate}
	\item Navigieren Sie zum Menü \menu{Werkzeuge}
	\item Wählen Sie den Unterpunkt \menu{Port} aus
	\item Selektieren Sie den verfügbaren Port, der mit Ihrem Arduino Nano 33 BLE Sense verbunden ist, wie in \ref{fig:arduinoPort}
\end{enumerate}

\begin{figure}[H]
	\centering
	\caption{Selektrieren des verfügbaren Ports}
	\label{fig:arduinoPort}
	\includegraphics[width=0.8\textwidth]{Arduino/ArduinoIDE/arduinoPort.png}
\end{figure}


\subsection{Sketch verifizieren und hochladen}

Nach der Auswahl des korrekten Ports folgt der Upload des Arduino-Programms. Gemäß Best Practices sollte das Programm \textbf{zuerst verifiziert} werden:

\begin{enumerate}
	\item Klicken Sie auf \menu{Überprüfen}
	\item Die IDE prüft nun auf Syntaxfehler und gibt Warnungen aus
	\item Bei erfolgreicher Verifizierung klicken Sie auf \menu{Hochladen} 
\end{enumerate}

\noindent\textbf{Verbindungsstatus:}
\begin{itemize}
	\item Der korrekte Verbindungsstatus wird in der unteren rechten Ecke der IDE angezeigt
	\item Dort sollten Board-Typ (Arduino Nano 33 BLE Sense) und Port aufgeführt sein
\end{itemize}

\noindent\textbf{Upload-Prozess:}
\begin{itemize}
	\item Der Compilierstatus wird im schwarzen Ausgabefenster angezeigt
	\item Nach erfolgreichem Upload blinkt die Onboard-LED (Pin 13) kurz auf
	\item Bei Port-Problemen: Port erneut unter \menu{Werkzeuge > Port} auswählen
\end{itemize}

\section{Konfiguration testen}

Zur Überprüfung der korrekten Einrichtung führen wir den Beispiel-Sketch \menu{Blink} aus:

\begin{enumerate}
	\item Öffnen Sie \menu{Datei > Beispiele > 01.Basics > Blink}
	\item Klicken Sie auf \menu{Hochladen}
	\item Die orangefarbene Onboard-LED (Pin 13) sollte nun im 1-Sekunden-Takt blinken
\end{enumerate}


\subsection{Test der Einrichtung mit dem Beispiel-Sketch \FILE{blink.ino}}

Die Arduino IDE bietet zahlreiche Beispiele zum Testen der Hardware. Zur Überprüfung der korrekten Konfiguration eignet sich besonders das grundlegende Beispiel \FILE{blink.ino}, das nach dem Upload die Onboard-LED mit 1 Hz blinken lässt.

\noindent\textbf{Vorgehensweise:}

\begin{enumerate}
	\item \textbf{Beispiele öffnen:}
	\begin{itemize}
		\item Arduino IDE starten
		\item Im Menü \menu{Datei} auf \menu{Beispiele} zeigen
	\end{itemize}
	
	\item \textbf{Beispielsketch auswählen:}
	\begin{itemize}
		\item Kategorie \menu{01.Basics} auswählen
		\item Sketch \menu{Blink} öffnen (siehe Abbildung~\ref{fig:arduinoBlink})
	\end{itemize}
	
	\item \textbf{Blink-Sketch analysieren:}
	\begin{itemize}
		\item Der geöffnete Sketch erscheint in neuem Fenster (Abbildung~\ref{fig:arduinoBlink2})
		\item Der Code blinkt die LED an Pin 13 im Sekundentakt
	\end{itemize}
	
	\item \textbf{Sketch uploaden:}
	\begin{itemize}
		\item Auf \menu{Hochladen} klicken (Pfeil-Symbol)
		\item Nach erfolgreichem Upload blinkt die orange Onboard-LED
	\end{itemize}
	
\end{enumerate}

\begin{figure}[H]
	\centering
	\includegraphics[width=0.8\textwidth]{Arduino/ArduinoIDE/arduinoBlink.png}
	\caption{Navigation zum Blink-Beispiel in der Arduino IDE}
	\label{fig:arduinoBlink}
\end{figure}

\begin{figure}[H]
	\centering
	\includegraphics[width=0.8\textwidth]{Arduino/ArduinoIDE/arduinoBlink2.png}
	\caption{Code des Blink-Beispiels in der Arduino IDE}
	\label{fig:arduinoBlink2}
\end{figure}

\noindent\textbf{Erfolgskontrolle:}
\begin{itemize}
	\item Die orange LED (Pin 13) blinkt im 1-Sekunden-Takt
	\item Im Ausgabefenster erscheint \glqq Upload abgeschlossen\grqq
	\item Keine Fehlermeldungen im schwarzen Konsolenfenster
\end{itemize}

Damit ist die Funktionsfähigkeit des Arduino Nano 33 BLE Sense bestätigt und das Board kann für anspruchsvollere Projekte verwendet werden



